\documentclass[12pt]{article}
\usepackage[utf8]{inputenc}
\usepackage[T1]{fontenc}
\usepackage{hyperref}
\title{Wadikali Language Resources for K-12 Curriculum Development}
\author{Compiled by Claude Sonnet 5 (claude-sonnet-5) for Andr\'as Kornai}
\date{}
\begin{document}
\maketitle

\section{Introduction}
Wadikali (ISO 639-3 wdk) is an extinct Australian language of the Yarli group, a close relative of Malyangapa and Yardliyawara. It was spoken in the north-east of South Australia, near the New South Wales and Queensland borders (sources differ on the location). Search results say it has not been spoken since the 1930s. Descendants of the Wadikali people still exist. It is known from a 72-word vocabulary in Tindale's 1934 notebook, recorded from Ned Palpilina.

\section{Primary Educational Resources}
None found.

\section{Linguistic and Reference Materials}
Glottolog entry: \url{https://glottolog.org/resource/languoid/iso/wdk}. The Yarli overview is at \url{https://en.wikipedia.org/wiki/Yarli_language}. Work by Hercus and Austin on the Yarli languages is mentioned in search results; one possible copy is at \url{https://www.academia.edu/2099279}, but I did not confirm its content.

\section{Religious and Community Resources}
None found.

\section{Digital and Audio Resources}
No recordings exist, since the language is not documented from speakers beyond the short word list.

\section{Curriculum Development Context}
The documentation is minimal. A heritage curriculum could cover the history and people of the area, with the Wadikali word list as a small component. Fuller revitalization work would have to draw on better-documented related Yarli languages, with community involvement.
\end{document}